% Revize 2026-10-28: obsah, matematika a návaznost; podrobnosti a časy
% viz planovani/REVIZE_PRIPRAV_1_ROCNIK.md.
% Aktuální porovnání se skutečným obsahem místních PDF Realisticky.cz:
% externi_zdroje/realisticky/01_Zakladni_poznatky/08_Mnohocleny/010810_Mnohocleny_shrnuti/010810_Lekce.pdf
% externi_zdroje/realisticky/01_Zakladni_poznatky/09_Lomene_vyrazy/010901_Kraceni_a_rozsirovani_lomenych_vyrazu/010901_Lekce.pdf
% externi_zdroje/realisticky/02_Funkce_a_rovnice/01_Linearni_funkce/020104_Funkce_definicni_obor/020104_Lekce.pdf
% externi_zdroje/realisticky/02_Funkce_a_rovnice/01_Linearni_funkce/020105_Graf_funkce_I/020105_Lekce.pdf
% Hodina 69, 1. ročník SPŠ. Rozsah: Rozpis_vyuky_MAT_1_rocnik.tex, řádek 69.

% Plán: 1_rocnik/tematicke_plany/SPS.pdf; společná projektová šablona.

% Učebnice: ucebnice/01_Zakladni_poznatky/Algebraicke_vyrazy/Algebraické výrazy.tex

% Původ úloh: Vlastní metodika a zadání k následné nápravě. Podklad funkcí: ucebnice/02_Funkce_a_rovnice/Zaklady_funkci/Funkce.tex, pojem funkce a práce s předpisem. Obsah se řídí řádkem 69 rozpisu.

% Metodický zdroj: externi_zdroje/realisticky/01_Zakladni_poznatky/09_Lomene_vyrazy/010901_Kraceni_a_rozsirovani_lomenych_vyrazu/010901_Lekce.pdf

% URL: http://www.realisticky.cz/ucebnice/01%20Matematika%20S%C5%A0/01%20Z%C3%A1kladn%C3%AD%20poznatky/09%20Lomen%C3%A9%20v%C3%BDrazy/01%20Kr%C3%A1cen%C3%AD%20a%20roz%C5%A1i%C5%99ov%C3%A1n%C3%AD%20lomen%C3%BDch%20v%C3%BDraz%C5%AF.pdf

% Metodický zdroj: externi_zdroje/realisticky/01_Zakladni_poznatky/08_Mnohocleny/010802_Nasobeni_mnohoclenu/010802_Lekce.pdf

% URL: http://www.realisticky.cz/ucebnice/01%20Matematika%20S%C5%A0/01%20Z%C3%A1kladn%C3%AD%20poznatky/08%20Mnoho%C4%8Dleny/02%20N%C3%A1soben%C3%AD%20mnoho%C4%8Dlen%C5%AF.pdf

\documentclass[11pt,a4paper]{article}

% Spolecny styl pro ucitelske pripravy MAT (Work/MAT)
% Verze 1.1, 2026-10-23
\NeedsTeXFormat{LaTeX2e}
\ProvidesPackage{priprava}[2026/10/23 v1.1 Sablona ucitelskych priprav MAT]

% Zakladni typografie a cestina
\RequirePackage[T1]{fontenc}
\RequirePackage[utf8]{inputenc}
\RequirePackage[czech]{babel}
\RequirePackage{lmodern}
\RequirePackage{microtype}
\RequirePackage{amsmath,amssymb}
\RequirePackage{xcolor}
\RequirePackage{enumitem}
\RequirePackage[most]{tcolorbox}
\RequirePackage{tabularx,booktabs,array}
\RequirePackage{fancyhdr}
\RequirePackage{needspace}
\RequirePackage{geometry}
\RequirePackage[hidelinks]{hyperref}

\geometry{a4paper,top=17mm,bottom=18mm,left=23mm,right=23mm}
\definecolor{notegray}{gray}{0.39}
\setlength{\parindent}{0pt}
\setlength{\parskip}{2.5pt plus 1pt minus 0.5pt}
\setlist[itemize]{topsep=2pt,itemsep=1pt,parsep=0pt,leftmargin=1.45em}
\setlist[enumerate]{topsep=2pt,itemsep=2pt,parsep=0pt,leftmargin=1.65em}
\setcounter{secnumdepth}{0}
\renewcommand{\arraystretch}{1.12}

% Zahlavi: vlevo rocnik | globalni tema, vpravo cislo: nazev hodiny
% Zapati zustava zcela prazdne.
\pagestyle{fancy}
\fancyhf{}
\renewcommand{\headrulewidth}{0.35pt}
\renewcommand{\footrulewidth}{0pt}
\setlength{\headheight}{14pt}
\newcommand{\nastavPripravu}[4]{%
  \fancyhead[L]{\small #1\enspace\textbar\enspace #2}%
  \fancyhead[R]{\small #3: #4}%
}

% Minimalisticke cernobile boxy. Kazda definice zvlast.
\tcbset{commonlesson/.style={enhanced,breakable=false,colback=white,colframe=black,
  colbacktitle=white,coltitle=black,boxrule=0.45pt,
  arc=0.8mm,outer arc=0.8mm,left=3mm,right=3mm,top=1.6mm,bottom=1.6mm,
  toptitle=1mm,bottomtitle=0.7mm,fonttitle=\bfseries,
  attach boxed title to top left={xshift=2.5mm,yshift=-1.15mm},
  boxed title style={colback=white,colframe=white,boxrule=0pt,left=1mm,right=1mm}}}
\newtcolorbox{zadaniBox}[1][]{commonlesson,title={Zadání},#1}
\newtcolorbox{prikladBox}[1][]{commonlesson,title={Příklad},#1}
\newtcolorbox{definiceBox}[1][]{commonlesson,title={Definice},#1}

% Spolecny box pro zadani a jeho reseni, bez mezery mezi castmi.
% Pouziti:
%   \begin{prikladReseniBox}
%     Zadani ulohy.
%     \tcblower
%     \begin{align*} ... \end{align*}
%   \end{prikladReseniBox}
% Spodni cast ma jemny sedy podklad a delici linku; lze v ni sazet
% vice radku matematiky. Obe casti maji jen jeden spolecny ramecek.
\tcbset{solutionlesson/.style={
  commonlesson,
  enhanced jigsaw,
  breakable,
  colbacklower=black!3,
  segmentation style={solid,draw=black!25,line width=0.35pt},
  fontlower=\small\color{notegray},
  before lower={\noindent\textbf{Řešení:}\enspace},
  leftlower=3mm,rightlower=3mm
}}
\newtcolorbox{prikladReseniBox}[1][]{solutionlesson,title={Příklad},#1}
\newtcolorbox{zadaniReseniBox}[1][]{solutionlesson,title={Zadání},#1}

% Pomucky pro autora pripravy: sedive poznamky; \reseni pro starsi pripravy
\newcommand{\poznamka}[1]{\par\begingroup\small\color{notegray}\itshape
  \textnormal{\textbf{Poznámka pro učitele.}}\enspace #1\par\endgroup}
\newcommand{\reseni}[1]{\par\begingroup\small\color{notegray}
  \textbf{Očekávané řešení:}\enspace #1\par\endgroup}
\newcommand{\blok}[1]{\needspace{5\baselineskip}\subsection*{#1}}
\newcommand{\vyrok}[1]{\uv{#1}}


\nastavPripravu{1. ročník}{Algebraické výrazy}{69}{Pololetní písemná práce}

\begin{document}

\blok{Výklad a souvislosti}

Rozpis vymezuje pololetní práci na algebraické výrazy a základy funkcí. Tato společná souhrnná hodina navazuje na algebraický blok 41–58 a na výuku funkcí v hodinách 59–68. Nepřidáváme novou hodinu ani nový obsah; tato příprava obsahuje organizaci a následnou nápravu, nikoli žákovské zadání testu.

\begin{definiceBox}[title={Připomenutí definice},breakable=false]
\textbf{Funkce} přiřazuje každému vstupu ze svého definičního oboru právě jednu funkční hodnotu. Bod $[x,y]$ leží na grafu právě tehdy, když $x$ patří do definičního oboru a $y=f(x)$.
\end{definiceBox}


Před hodinou připrav schválenou souhrnnou práci podle skillu \texttt{vytvor-test}. Ověř návaznost na skutečně probrané učivo do hodiny 68: úpravy a rozklady mnohočlenů, lomené výrazy s podmínkami, pojem funkce, čtení předpisu a grafu, lineární a konstantní funkci. Grafy s absolutní hodnotou začínají až v hodině 70 a do této práce je nezařazujeme.

Převážnou část hodiny ponech samostatné písemné práci; přesnou organizaci a povolené pomůcky stanoví učitel u konkrétního zadání. Tato příprava další pomůcky nepovoluje. Následující příklady používej až po odevzdání nebo při individuální zpětné vazbě po opravě; pokud nezbude čas, jsou zásobou pro konzultaci.

\blok{Řešené příklady}

\begin{prikladReseniBox}[title={Příklad 1},breakable=false]
Při rozboru uprav $\frac{x^2-4}{x-2}$ a vysvětli zachování podmínek.
\tcblower
Pro $x\ne2$ lze krátit $x-2$ v rozkladu $(x-2)(x+2)$, takže dostaneme $x+2$. Hodnota $x=2$ zůstává vyloučena, ačkoli upravený zápis má pro ni číselnou hodnotu.
\end{prikladReseniBox}


\begin{prikladReseniBox}[title={Příklad 2},breakable=false]
Při následné nápravě pro $f(x)=2x+3$, $x\in\mathbb R$, určete $f(-1)$, $f(0)$, $f(1)$ a popište sestrojení grafu.
\tcblower
Hodnoty jsou $1$, $3$ a $5$. Graf je přímka určená například body $[0,3]$ a $[1,5]$; bod $[-1,1]$ poskytne další kontrolu. Zvýšení argumentu o $1$ zvýší funkční hodnotu o $2$.
\end{prikladReseniBox}


\blok{Samostatné procvičení}

\begin{zadaniReseniBox}[title={Procvičení},breakable=false]
\begin{enumerate}[label=\alph*)]
\item Po opravě rozlož $x^2-9$ a výsledek ověř.
\item Pro následné procvičení určete hodnotu $h(2)$ u $h(x)=-x+4$ na $\mathbb R$ a dva body jejího grafu.
\item Při rozboru rozhodni, zda body $[0,3]$ a $[2,3]$ leží na grafu $f(x)=2x+3$ na $\mathbb R$.
\end{enumerate}
\tcblower
\begin{enumerate}[label=\alph*)]
\item $(x-3)(x+3)$; roznásobením vznikne $x^2-9$ pro všechna reálná $x$.
\item $h(2)=2$; například body $[0,4]$ a $[4,0]$ leží na grafu, protože splňují předpis.
\item První ano: $f(0)=3$. Druhý ne: $f(2)=7\ne3$. Nestačí shodná jedna souřadnice.
\end{enumerate}
\end{zadaniReseniBox}


\poznamka{Rozlišuj záměnu argumentu a funkční hodnoty, chybu při dosazení a chybu v samotné algebraické úpravě. Výsledky jednotlivců nerozebírej veřejně. Konkrétní test, jeho varianty a hodnocení se připravují samostatně; zde nejsou body ani známkové hranice.}

\end{document}
